\subsection{The original two links and their common Hecke ideal}
\label{subsec:rh-rank-four-link-ideal}

We retain the fitted deep tower $T$ of the original rank-four problem.
Write $F=T.F$, $K=T.K$, $L=FK$, and $S=D.S$ for the original finite
set of rational primes. Thus $F$ is totally real abelian over
$\mathbb Q$, $K$ is imaginary quadratic, $L/F$ is unramified at all
finite places, and the primes of $S$ (including $2,\ell$) split in $K$.
Let $\epsilon$ be the $\ell$-adic cyclotomic character,
$\omega$ its Teichmuller part, and $e=D.e$. The source characters are
\[
 \alpha=\epsilon\omega^{-1}\xi,\qquad \Psi=\epsilon^e\psi,\qquad
 \nu=\alpha\Psi .
\]
Here $\xi$ is finite of order prime to $\ell$, $\psi$ is finite,
$\xi(c)=-1$, $\Psi(c)=1$, and both are unramified outside $S$.
The four characters are $(\nu,\Psi,\alpha,1)$ in the original
ordering, or its prescribed relabeling. The deep condition makes
$\psi|_{G_F}$ of order prime to $\ell$. No new tower extension is
made in this subsection.

\begin{lemma}[Characters for the original link]
\label{lem:rh-four-link-characters}
After one finite enlargement of the coefficient field, independent
of $X$, there are characters $\kappa,\beta:G_L\to\mathcal O^\times$
and $t:G_L\to\mathbb Z_\ell$ satisfying the character conditions
of the original \texttt{Link}, including the odd squared inertial
exponent for $\beta$ and the integral inertial exponent for $\kappa$.
\end{lemma}
\begin{proof}
Let $\psi_0$ be the prime-to-$\ell$ part of $\psi$, and set
$\eta=(\omega^{e+3}\xi\psi_0)^{-1}$. This is a finite character
of $G_{\mathbb Q}$ of order prime to $\ell$, unramified outside $S$.
It is even: $\omega(c)=-1$, $\psi(c)=(-1)^e$, and passage to the
prime-to-$\ell$ part does not change this value. On $G_F$,
$\psi_0=\psi$.

We first construct a finite character $\kappa_0$ of $G_K$ of order
prime to $\ell$, unramified outside $S$, with
$\kappa_0\kappa_0^c=\eta|_{G_K}$. Use global reciprocity. Choose
a symmetric finite ray modulus of $K$ supported above $S$,
with exponents large enough to contain the conductor of the
rational idele character corresponding to $\eta$ at each split
component. On the image of the diagonal rational finite unit
ideles in the finite ray class group prescribe $\eta$. This
prescription is well defined: an element of that diagonal unit
group which is trivial in the ray class group agrees modulo the
ray unit subgroup with a global unit of $K$. The units of $K$
are $\{\pm1\}$, since $2$ splits in $K$ (the two exceptional
imaginary quadratic unit fields do not have this property).
The evenness of $\eta$ and the conductor choice therefore kill
the kernel. Every rational finite idele can be made a unit idele
by multiplication by a rational number. The finite-group character
extends to the whole ray class group since the roots of unity
in the algebraic closure are divisible. Taking its prime-to-$\ell$
part leaves the prescribed restriction unchanged. This gives
$\kappa_0$. For an idele $z$ of $K$, the product $zz^c$ is the
diagonal image of its norm. Compatibility of reciprocity with the
norm proves $\kappa_0\kappa_0^c=\eta|_{G_K}$, with the same
reciprocity convention on the two fields.

Next use the maximal abelian pro-$\ell$ quotient of $G_K$ unramified
outside $\ell$. Its local principal-unit subgroup has finite index
by the ray class exact sequence, and is a product of two rank-one
$\mathbb Z_\ell$-modules because $\ell$ splits in $K$ and its global
units are finite of order prime to $\ell$. Complex conjugation
interchanges the two factors. The torsion-free anticyclotomic
quotient has rank one; identify it with $\mathbb Z_\ell$ and denote
the resulting character by $t_K$. Both local inertia maps into
this quotient are nonzero, have open images, and are negatives
under conjugation. Choose one of the two places and identify the
cyclotomic principal-unit image on its inertia with
$1+\ell\mathbb Z_\ell$. By local reciprocity, the homomorphism
$t_K$ on this pro-$\ell$ part is a nonzero homomorphism of
rank-one $\mathbb Z_\ell$-modules. If its image is
$\ell^b\mathbb Z_\ell$, choose an inertia element mapping to
$\ell^b$ and write $v$ for its cyclotomic principal-unit value.
Choose a root $u\equiv1$ modulo the maximal ideal with
$u^{2\ell^b}=v$ in a finite coefficient extension. Such a root
exists: first take an $\ell^b$-th root of $v$ (its reduction is
$1$), then its unique principal-unit square root, using odd $\ell$.
Put $\beta_K(g)=u^{t_K(g)}$ by the continuous binomial power.
The defining equality on a topological generator gives
$\beta_K^2=\epsilon/\omega$ on the principal-unit part of that
inertia group. At the conjugate place the exponent is $-1$.
Thus on inertia $\beta_K^2\epsilon^{-b_v}$ has finite image
with $b_v=1$ or $-1$ respectively. In particular the required
exponent is odd. The character is anticyclotomic, principal-unit
valued, and trivial away from $\ell$ on inertia.

Put $\gamma=(\epsilon/\omega)^{1/2}$, the unique principal-unit
square root, and define the characters on $G_L$ by restriction
and the formula
\[
 \beta=\beta_K|_{G_L},\qquad t=t_K|_{G_L},\qquad
 \kappa=\kappa_0|_{G_L}\,\gamma^{-(e+4)}\beta^{e\bmod 2}.
\]
These are the source's characters, with
\[
              \kappa\kappa^c\,\nu\,\epsilon^3=1,\qquad
              \beta\beta^c=1,\qquad t^c=-t .
\]
Indeed $\psi=\psi_0$ on $G_L$ and
$\gamma^2=\epsilon/\omega$, so the first equality is a direct
multiplication. The characters are continuous and take values
in a fixed finite coefficient field. Outside $S$, $\kappa$ is
unramified; on inertia away from $\ell$ its image has order
prime to $\ell$ because only the finite factor $\kappa_0$ can
contribute. At each $\ell$-adic place, up to finite order on
inertia, the exponent of $\kappa$ is
$[-(e+4)+(e\bmod2)b_v]/2$, an integer.
If a local element has an $\ell$-power in the closure of the
commutator subgroup, $\kappa_0$ kills it by its prime-to-$\ell$
order, $t$ kills it by torsion-freeness, and $\epsilon/\omega$
kills it by torsion-freeness of $1+\ell\mathbb Z_\ell$.
Consequently $\kappa$ and $t$ kill the prescribed local
$\ell$-primary torsion subgroups. This verifies the local
torsion clause of the original link (also after passage to
its unramified-outside-$S'$ quotient). The condition
$\ker t\subset\ker\beta$ is immediate.
\end{proof}

\begin{theorem}[Two original links and determinant maps]
\label{thm:rh-four-link-maps}
For every fitted deep tower the two links in
\texttt{exists\_two\_links\_same\_kappa\_not\_dvd\_of\_fit\_deep}
exist with the same $\kappa,\beta,t$ and distinct distinguished
rational primes. For every sufficiently large original coefficient
extension and every original space $X$, the original determinant
maps $\phi_n:B_0\to X.B$ exist for each $n\geq0$ and each link.
\end{theorem}
\begin{proof}
The finite residual diagonal characters on $G_L$ have the
source's opposite-pair product relation, so their quotients
generate the group of characters generated by the two root
coordinates $w\in\mathbb Z^2$. The fitted hypothesis
\texttt{Ok.g2} says that $\bar\omega$ is not in this group.
For the finite abelian joint image of the characters, elementary
duality with characters into the finite coefficient field's
multiplicative group identifies the double annihilator of any
subgroup with that subgroup: the exponent of the joint image
divides the order of that cyclic multiplicative group.
There is therefore an element of $G_L$ on which all four
residual diagonal characters agree and $\bar\omega\ne1$.

Apply Chebotarev to a finite Galois extension of $\mathbb Q$
containing $L$ and the fields cut out by the residual diagonal
characters and $\bar\omega$. Choose two distinct rational primes
outside $S$, the discriminant of $L$, and $2$, with Frobenius
in the indicated conjugacy class. These primes split completely
in $L$ (as $L/\mathbb Q$ is Galois), and at either chosen place
of $F$ their residue norm is the rational prime. On its
decomposition group the four residual characters agree
and $\bar\epsilon\ne1$, since inertia is trivial and the
unramified group is generated topologically by Frobenius.
In particular $L$ contains no primitive $\ell$-th root of
unity. For each link take the set of places of $F$ over $S$
together with the respective single distinguished place.
They all split in $L/F$; at $\ell$ keep precisely the
original choices $T.\iota$. This is the Hida setting
\texttt{SetD} in the source. The original level with
$V.S'=s.S$ is supplied by the established ordinary
level construction. Equivalently its Galois quotient is the
kernel of the sign character in the maximal Galois group of
$F$ unramified outside $s.S$; the quadratic extension
$L/F$ is itself unramified at every finite place, so that
kernel is the corresponding quotient for $L$.

For both levels use the characters from
Lemma~\ref{lem:rh-four-link-characters}. The four residual
characters are exactly $\overline{\kappa\chi_i}|_{G_L}$,
distinct by \texttt{Ok.g3}; their ordered local characters
are their restrictions, and their prescribed local torsion
characters are $1$. The conjugate pairing on determinants
follows from $\kappa\kappa^c\nu=\epsilon^{-3}$ and the
opposite-pair product relation. The residual characters
factor through the level quotient, by the unramifiedness
outside $S$ of all factors. At the distinguished prime
$\kappa$ is unramified, so the preceding Frobenius choice
gives \texttt{Resid.V0ok}. The coefficient field can be
enlarged once to contain the characters, the images of
all embeddings of $L$ into $\overline{\mathbb Q}_\ell$,
and the original coefficient field. Every further coefficient
extension then has the required largeness. This constructs
the original two links without replacing either distinguished
prime by a prime of higher residue degree.

Write $B=X.B$ and $A_X=X.A$; the original space includes
the injective coefficient map $B\to A_X$, its universal
symplectic matrices $R(g)$, labeled flags, and the
coefficients of $C_g(Z)=\det(Z-R(g))$ and local diagonal
characters in $B$. Restrict the matrices to $G_L$ and twist
by $\kappa\beta^n$. All their characteristic polynomials
have coefficients in $B$, with exactly the prescribed root
scaling. Their determinant on the group algebra is also
defined over $B$: the integral Amitsur formula for
$\det(1-\sum_j Z_j M_j)$ expresses its coefficients using
the characteristic-polynomial coefficients of words in
the $M_j$. This formula is a universal polynomial identity
over $\mathbb Z$ (no division by $4!$); in this application
every word is again the matrix of a group element. By the
injectivity $B\to A_X$, the polynomial-law identities hold
over $B$ and after every scalar extension of $B$.
The determinant is continuous in the maximal-ideal topology
by the original coefficient-continuity conditions.

Here are the exact descent and local checks. Outside the
chosen set of primes, the original residual representation
and $\kappa$ are trivial on inertia, and the original
symplectic killing condition and unramifiedness of $\beta$
make the twisted matrices trivial there. The determinant
therefore descends to the level quotient: the matrices kill
the normal subgroup generated by these inertia groups;
continuity of the characteristic-polynomial coefficients
into the separated ring $B$ extends this invariance to
its closure. Amitsur's formula gives the same invariance
for the whole polynomial law. This argument does not
require $A_X$ itself to be complete. At an old inertia
place, if all four twisted residual characters are one,
the diagonal list includes the trivial original character,
so $\bar\kappa=1$ and all original residual characters
are one. The original killing condition on the residual
kernel makes $R$ the identity there; $\kappa$ is of
prime-to-$\ell$ order on that inertia, hence it too is
one, and $\beta$ is unramified. This gives exactly the
\texttt{Pair.I} trace relations. The ordinary product
relations hold by the ordered upper triangular flags at
the original local decomposition groups: the rightmost
factor kills the last quotient and the successive factors
lower the image through the flag. Multiplying by the
twisting character preserves these identities. The
prescribed local torsion characters are one by the original
\texttt{Fr.IsPair} conditions and the torsion check of the
lemma. Finally $R(c)$ conjugates $R(g)$ to $R(cgc^{-1})$,
and the symplectic relation and
$\kappa\kappa^c\nu\epsilon^3=1$ prove the determinant
polarization with multiplier $\epsilon^{-3}$. Thus the
twisted determinant and the twisted labeled local characters
constitute a literal \texttt{Pair} over $B$ for the
original residue datum. Its universal property gives
$\phi_n$ and exactly the equalities of \texttt{Link.IsPhi}.
\end{proof}

\begin{theorem}[The original image ideal: finite weights and seed]
\label{thm:rh-four-image-ideal}
For the two links of Theorem~\ref{thm:rh-four-link-maps}, let
$J_0,J'_0$ be their actual ordinary Hecke ideals and
$\phi_0,\phi'_0$ their untwisted maps. Form
\[
                   J=\phi_0(J_0)B+\phi'_0(J'_0)B .
\]
This ideal satisfies the original \texttt{fin} and
\texttt{seed} clauses of \texttt{PotentialModularityTower}.
For arithmetic weights $D.wts4=(e+1,e,1,0)$ with $e\geq2$,
it also satisfies the \texttt{pur} clause. At every such
point where the representation over $G_L$ is irreducible,
the untwisted representation over $G_L$ has a cuspidal
regular algebraic attachment.
\end{theorem}
\begin{proof}
For finiteness it suffices to use just the first image ideal.
Let $\mathfrak q$ be a prime containing that ideal, the
coefficient uniformizer, and all original inertial
differences $\psi_i(s)-1$. The domain $C=B/\mathfrak q$
is a complete local algebra of characteristic $\ell$ with
the same finite residue field $k$ as the coefficient DVR.
The composite $B_0\to C$ kills $J_0$. Its restriction to
the original weight ring is the residual evaluation:
the inertial characters of the link at $\ell$ are those
of $X$ multiplied by the fixed $\kappa$, which reduce
in $C$ to their prescribed residues. The universal
weight property, or the original weight generators,
gives the assertion for the whole weight ring.
Since $B_0/J_0$ is finite over the weight ring, its image
in $C$ consists of elements integral over $k$.

It follows that all coefficients of $C_g$ for $g\in G_L$
are algebraic over $k$ (the scalar twist is a fixed residue
unit), as are all local characters on the chosen ordinary
groups (these groups lie in $G_L$, since the places split).
For $g\in G_F$, $g^2\in G_L$. The characteristic-polynomial
coefficients of a matrix are integral over those of its square:
with universal eigenvalues $a_1,\ldots,a_4$, take the monic
polynomial in $Z$ with roots
$e_i(\pm a_1,\ldots,\pm a_4)$ for all sixteen sign choices,
where $e_i$ is the $i$-th elementary symmetric polynomial.
Its coefficients are symmetric polynomials in the squares
with integral coefficients. It gives a universal matrix
identity, as can be seen first for matrices with distinct
eigenvalues over characteristic zero, then by polynomial
identity over $\mathbb Z$. The original map $B\hookrightarrow A_X$
lets this identity descend to $B$, so it also holds in $C$.
This proves that every $C_g$ coefficient is algebraic over $k$.

An element of a local domain with residue field $k$ which
is algebraic over the finite field $k$ belongs to $k$:
subtract its residue; a nonzero algebraic element over $k$
would be a unit and cannot lie in the maximal ideal.
Thus every characteristic-polynomial coefficient and every
labeled local character equals its residual value.
The exact \texttt{Space.genB} uniqueness property now
makes $B\to C$ equal to residual evaluation. Consequently
$\mathfrak q=\mathfrak m_B$. Since $B$ is Noetherian,
the ideal in the \texttt{fin} condition contains a power
of $\mathfrak m_B$. This proves the precise finite-weight
clause for $J$.

The seed-support theorem of the preceding subsection,
applied to each link and using the parity correction and
the twist-ideal equality (RS1), gives $J\subset\mathfrak q$
for every prime of $B$ below the kernel of any original
seed. This is exactly the \texttt{seed} clause.

For the last assertions take a characteristic-zero point
with the stated arithmetic weights and killing $J$.
The first link with $\kappa$ already has integral inertial
exponents: the local squared weights are twice the integers
$D.wts4$, and $\kappa$ has integral local exponents.
They are strictly ordered because $e\geq2$. The point is
therefore an arithmetic point of the actual ordinary
Hecke quotient. Apply its already established
arithmetic-point decomposition theorem. Each Speh
multiplicity greater than one would give two successive
constituents $\sigma,\sigma\otimes\epsilon^{-1}$ in the
semisimple Galois representation. Choose stable integral
lattices in these continuous characteristic-zero
constituents and reduce. The residual semisimplification
of the whole representation is the given multiplicity-free
sum of the four residual characters. Hence some two
residual characters differ by $\bar\omega$ (take any
constituent of the first reduced lattice and its twist).
The common factor $\bar\kappa$ cancels in the quotient.
This contradicts \texttt{Ok.g2}. Thus all Speh
multiplicities are one. The ordinary-family purity theorem
now gives equal complex absolute values of the good
Frobenius eigenvalues, in its prescribed normalization.
Undoing the common scalar twist preserves the equality.
For the rational representation in the \texttt{pur} clause
restrict Frobenius to the open subgroup $G_L$:
at primes unramified in the fields and representations,
Frobenius elements in this subgroup are powers of the
corresponding Frobenius elements over $L$. The equality of
absolute values persists under positive powers. This
implies the strict ratio bound of \texttt{PureW}, with
the open subgroup and the finite bad-prime product
required by that definition.

If the restricted representation is absolutely irreducible,
the established ordinary arithmetic classicality theorem
gives a cuspidal regular algebraic conjugate-self-dual
attachment to the twisted representation. Untwist by
$\kappa^{-1}$, which is an algebraic Hecke character by
the integral local exponent and finite tame inertia
conditions. This proves the asserted untwisted
cuspidal attachment over the actual $L$.
\end{proof}

The last conclusion is deliberately weaker than the original
\texttt{aut} clause of \texttt{PMExAlt}: it does not assert
the simultaneous strict local windows in \texttt{WitA}, or
their descent to every imaginary quadratic field. The
characteristic-prime \texttt{p1} clause is also separate.
In particular neither is a consequence solely of the two
image-ideal construction and characteristic-zero seed
support proved here.
